
\section{Limitations and Future Work}
First, due to the computational resources limit, the empirical evaluation is concentrated on mathematical reasoning, leaving its generalizability to other complex reasoning domains (e.g., commonsense or strategic planning) an open question. Second, even on these tasks, while STARS prevents catastrophic collapse and enables scaling beyond the training horizon, performance does not always improve monotonically with more steps. This indicates that achieving fully reliable and predictable test-time scaling remains a challenge for the most difficult problems.

Future work will explore extending the STARS framework to a wider variety of reasoning and planning tasks, and investigating adaptive mechanisms to more robustly guide latent trajectories toward optimal, stable fixed points.
\section{Training Algorithm Details}
% \begin{algorithm}[h]
%    \caption{STARS Training for Looped Language Models}
%    \label{alg:jsrr}
% \begin{algorithmic}
%    \STATE {\bfseries Input:} Dataset $\mathcal{D}$, Initial parameters $\theta$, Distribution $\mathcal{P}$, Regularization weight $\lambda$, Learning rate $\eta$
%    \REPEAT
%    \STATE Sample a batch $\{\mathbf{x}^{(i)}\}_{i=1}^N \sim \mathcal{D}$
%    \STATE Sample loop depth $t \sim \mathcal{P}$
%    \FOR{$i = 1$ $\mathbf{to}$ $N$}
%    \STATE \textit{// Forward pass to $t$}
%    \STATE $\mathbf{h}^{(0)}_i \gets \text{prelude}(\mathbf{x}^{(i)})$
%    \STATE $\mathbf{h}^{(t)}_i \gets \underbrace{\Phi_\theta \circ \dots \circ \Phi_\theta}_{t \text{ times}} (\mathbf{h}^{(0)}_i)$
   
%    \STATE \textit{// One step power iteration}
%    \STATE $\mathbf{v}^{(i)} \sim \mathcal{N}(0, I)$
%    \STATE $\mathbf{v}^{(i)} \gets \mathbf{v}^{(i)} / \|\mathbf{v}^{(i)}\|_2$
%    \STATE $\mathbf{u}^{(i)} \gets \text{JVP}( \mathbf{h}^{(t)}_i, \mathbf{v}^{(i)})$
%    \STATE $\mathbf{u}^{(i)} \gets \mathbf{u}^{(i)} / \|\mathbf{u}^{(i)}\|_2$
%    \STATE \textit{ // Estimate spectral radius via the Rayleigh
% quotient}
%    \STATE $\rho_t^{(i)} \gets \frac{\mathbf{u}^{(i)\top} \text{JVP}( \mathbf{h}^{(t)}_i, \mathbf{u}^{(i)})}{\mathbf{u}^{(i)\top} \mathbf{u}^{(i)}}$
%    \ENDFOR
   
%    \STATE \textit{// Loss computation following Equation \ref{e1} and \ref{e2}}
%    % \STATE $\mathcal{L}_{\text{SFT}}^{(t)} \gets \frac{1}{N} \sum_{i=1}^N -\log p_{\theta}^{(t)}(x_{\ell+1}^{(i)} \mid x_{1:\ell}^{(i)})$
%    % \STATE $\mathcal{L}_{\text{JSRR}}^{(t)} \gets \frac{1}{N} \sum_{i=1}^N \rho_t^{(i)}$
%    \STATE $\mathcal{L} \gets (1 - \lambda)\mathcal{L}_{\text{SFT}}^{(t)} + \lambda \mathcal{L}_{\text{JSRR}}^{(t)}$
   
%    \STATE \textit{// Optimization step}
%    \STATE $\theta \gets \theta - \eta \nabla_\theta \mathcal{L}$
%    \UNTIL{convergence}
% \end{algorithmic}
% \end{algorithm}
\begin{algorithm}[h]
   \caption{STARS Training for Looped Language Models with JSRR}
   \label{alg:stars_jacobian}
\begin{algorithmic}
   \STATE {\bfseries Input:} Dataset $\mathcal{D}$, Initial parameters $\theta$, Distribution $\mathcal{P}$, Regularization weight $\lambda$, Learning rate $\eta$, Power iteration steps $K$
   \REPEAT
   \STATE Sample a batch $\{\mathbf{x}^{(i)}, \mathbf{y}^{(i)}\}_{i=1}^N \sim \mathcal{D}$
   \STATE Sample loop depth $t \sim \mathcal{P}$

   \FOR{$i = 1$ $\mathbf{to}$ $N$}
      \STATE \textit{// Forward pass to loop depth $t$}
      \STATE $\mathbf{h}^{(0)}_i \gets \text{prelude}(\mathbf{x}^{(i)})$
      \STATE $\mathbf{h}^{(t)}_i \gets \underbrace{\Phi_\theta \circ \dots \circ \Phi_\theta}_{t \text{ times}}(\mathbf{h}^{(0)}_i)$

      \STATE \textit{// Initialize random direction}
      \STATE $\mathbf{v}^{(i)} \sim \mathcal{N}(0, I)$
      \STATE $\mathbf{v}^{(i)} \gets \mathbf{v}^{(i)} / \left(\|\mathbf{v}^{(i)}\|_2 + \epsilon\right)$

      \STATE \textit{// Power iteration using Jacobian-vector products}
      \FOR{$k = 1$ $\mathbf{to}$ $K$}
         \STATE $\mathbf{j}^{(i)} \gets \text{JVP}\left(\Phi_\theta, \mathbf{h}^{(t)}_i, \mathbf{v}^{(i)}\right)$
         \STATE $\mathbf{v}^{(i)} \gets \mathbf{j}^{(i)} / \left(\|\mathbf{j}^{(i)}\|_2 + \epsilon\right)$
      \ENDFOR
      \STATE $\mathcal{L}^{(i)}_{\text{JSRR}} \gets \|\mathbf{j}^{(i)}\|_2^2$
   \ENDFOR

   \STATE \textit{// Loss computation}
   \STATE $\mathcal{L}_{\text{SFT}}^{(t)} \gets \frac{1}{N} \sum_{i=1}^N -\log p_{\theta}^{(t)}(\mathbf{y}^{(i)} \mid \mathbf{x}^{(i)})$
   \STATE $\mathcal{L}_{\text{JSRR}}^{(t)} \gets \frac{1}{N} \sum_{i=1}^N \mathcal{L}^{(i)}_{\text{JSRR}}$
   \STATE $\mathcal{L} \gets \mathcal{L}_{\text{SFT}}^{(t)} + \lambda \mathcal{L}_{\text{JSRR}}^{(t)}$

   \STATE \textit{// Optimization step}
   \STATE $\theta \gets \theta - \eta \nabla_\theta \mathcal{L}$

   \UNTIL{convergence}
\end{algorithmic}
\end{algorithm}

\section{More Results}
\label{more_res}

% \subsection{Analysis}
\paragraph{Random loop sampling on PreNorm and PostNorm}
We supplement the experimental results of random loop sampling for PreNorm with LN and PostNorm with LN under three distributions (Log-Normal, Poisson, and Uniform), each with two configurations, as shown in Figure \ref{app:analysis} and Table \ref{tab:loop_params_app}.
\label{app:analysis}
\begin{figure}[htbp]
    \caption{The results with the random loop strategy across distributions and parameter sets for PreNorm with LN and PostNorm with LN.}
	\centering
	\includegraphics[width=0.95\linewidth]{fig/pre_post_random.jpg}
\end{figure}


\begin{table}[hbp]
    \centering
    % \small
    \caption{Hyperparameter settings for random loop distributions with PreNorm and PostNorm. Range indicates the clipping bounds $[\text{min}, \text{max}]$.}
    \label{tab:loop_params_app}
    % \setlength{\tabcolsep}{5pt} % 调整列间距
    \resizebox{0.5\textwidth}{!}{ 
    \begin{tabular}{l c c}
        \toprule
        \textbf{Distribution} & \textbf{Set 1 Configuration} & \textbf{Set 2 Configuration} \\
        \midrule
        {Log-Normal}                 & $\mu=2.62, \sigma=0.60$ & $\mu=3.2, \sigma=0.45$ \\
                                    & range: $[1, 40]$ & range: $[1, 80]$ \\
        \midrule
        {Poisson}                    & $\lambda=5$ & $\lambda=10$ \\
                                    & range: $[1, 30]$ & range: $[1, 30]$ \\
        \midrule
        Uniform                     & range: $[1, 10]$ & range: $[1, 30]$ \\
        \bottomrule
    \end{tabular}}
\end{table}


\paragraph{Hyperparameter analysis.}
We conducted a hyperparameter analysis by evaluating the regularization weight $\lambda$ on multi-digit addition tasks, with the results illustrated in \ref{fig:lambda_anal}. This investigation involved plotting the performance curve and the PCA-projected trajectory for a range of $\lambda$ values. The weight $\lambda$ directly influences the strength of the JSRR constraint, which, in turn, is reflected in the resulting latent dynamics. While our method demonstrates that different weights lead to the model converging to attractor, as shown in the trajectory plots, if the weight $\lambda$ is too large (e.g. 0.15 and 0.20), the model struggles to learn to solve the original multi-digit task. This is evident in the corresponding performance curves, which show a significant drop in accuracy. Therefore, although our method exhibits a degree of robustness, the regularization weight should not be excessively large, as it can impede the model's effective learning capabilities.
\begin{figure}[H]
	\centering
	\includegraphics[width=0.9\linewidth]{fig/weight.jpg}
    \caption{Hyperparameter analysis on the multi-digits addition task.}
     \label{fig:lambda_anal}
\end{figure}

\paragraph{Efficiency analysis during training phase.}
We conducted an efficiency analysis during the training phase for these four types, with Ouro-SFT as the baseline. The results are shown in Table \ref{tab:efficiency_analysis}. From the table, it can be observed that the efficiency of Ouro-Random Loop, Ouro-JSRR, and Ouro-STARS during the training phase is 1.4976, 1.5317, and 2.0439, respectively, relative to Ouro-SFT.

\begin{table}[H]
    \centering
    \caption{Efficiency analysis during training phase. The values are presented relative to Ouro-SFT as the base.}
    \label{tab:efficiency_analysis}
    \begin{tabular}{l c c c c}
        \toprule
        \textbf{Model} & \textbf{Ouro-SFT} & \textbf{Ouro-Random Loop} & \textbf{Ouro-JSRR} & \textbf{Ouro-STARS} \\
        \midrule
        \textbf{Relative Value} & 1 & 1.4976 & 1.5317 & 2.0439 \\
        \bottomrule
    \end{tabular}
\end{table}

\paragraph{Comparative Analysis and Ablation Study on AMC23.}
In addition to the mathematical reasoning benchmark analysis presented in the main text (Figure \ref{fig:ablation}), we further conducted a comparative analysis of our method against baselines and its ablation variants on the AMC23 dataset shown in Figure \ref{fig:amc23_ablation}.
\begin{figure}[H]
	\centering
	\includegraphics[width=\linewidth]{fig/amc23_ablation.png}
    \caption{Comparative analysis of our method against baselines and its ablation variants on AMC23.}
    \label{fig:amc23_ablation}
\end{figure}